\section{Introduction}
An experiment acquires reports by preparations and interventions with a causal order and a physical cost. Its executable future tests determine the predictive distinctions worth retaining. Geometry enters only after this experiment, its acquired law and its scored task have been fixed. The central question is therefore not merely how to compress a prescribed distribution. It is when an acquired geometric profile describes the optimal performance of a legal finite-resource causal observer, and how the observer and the converse can be obtained from primitive experimental data.

Three obstacles must be treated together. Selection pays for rejected reports and changes the law that is eventually used. A small unnormalized error can be amplified at a rare stopping event. A label used to route a live protocol may also serve as a terminal output when commitment is an event rather than an intrinsic state type. Omitting any of these distinctions can produce an incorrect resource curve. A further obstacle concerns quantifiers: one common controller for an unknown experiment is different from one optimized controller for each model or prior.

For a genuinely prepared comparison law, let $f(x)=\E_0[Z\mid x]$ and $w(x)=\E_0[\|Z-f(x)\|^2\mid x]$. Its squared-loss profile is
\[
 \Q_m(b)=\min_{|C|\le m,\ 0\le h\le b,\ \int h\,d\nu=1}
     \int [w+\dist(f,C)^2]h\,d\nu.
\]
The density cap records paid acquisition; $m$ records actual decoder width. This is a marked population trimmed-quantization functional. Population trimming and quantization have a classical statistical history \cite{CGM97,GG99}; the functional is not claimed as a new invention. Its causal interpretation requires an all-strategy lower certificate in one actual model and an all-model upper construction using the declared interface. The inherited Theorem~\ref{thm:transfer} proves that transfer on actual pre-call occupation, without running a competitor in a fictitious reference experiment where it might fail to terminate.

The present development strengthens that theory in three connected directions. First, Theorem~\ref{thm:width-defect} identifies the attained least error for the actual selected task at every finite decoder width. Its positive-envelope representation retains a single dictionary across reports and histories. In squared prediction, Theorem~\ref{thm:quadratic-defect} reduces the zero-defect criterion to the minimum of a scalar quadratic over the actual decision carrier. Full affine-square matrix order is sufficient but strictly stronger: two audit laws can agree on every squared prediction while their task matrices have incompatible kernels. The sharp example in Theorem~\ref{thm:sharp-width} has defect $\min(m,k)/(5k)$ in dimension $k$. Thus an unrestricted-width discrepancy can overcharge a fixed decoder by an unbounded factor. This hierarchy concerns positive error magnitude; its zero-defect order is already determined by localized one-value tests. It is not a new Blackwell hierarchy of zero-error experiment orders \cite{Blackwell,BR14}. Theorem~\ref{thm:width-transfer} carries these exact errors through the paid stopping proof and includes bounded decision-continuous and Bregman losses.

Second, we replace a search among specified one-call mixtures by two intrinsic synthesis results. Theorem~\ref{thm:causal-hierarchy} characterizes all proper full-history programs on a finite-branching classical or quantum interface with a positive call-cost floor. One common program attains a finite list of bounded terminal task moments at budget $N$ exactly when every finite-horizon problem for every finite set of operating models is feasible. The finite problems retain both uncommitted mass and the cost of its next mandatory call. Their sharp survival bound is $\min\{1,(N/\underline c-1)/T\}$, including the boundary $N=\underline c$. An infeasible infinite synthesis problem therefore has a finite obstruction and the strict dual certificate of Theorem~\ref{thm:synthesis-dual}. This is an exact common-policy existence theorem, not a finite test for arbitrary infinite data or a bounded-memory assertion. A legal fallback turns a feasible finite prefix into a paid approximation with explicit score and cost errors.

Theorem~\ref{thm:memory-certificate} answers the separate question at fixed classical memory for time-homogeneous finite-dimensional models. The continuing completely positive map determines a least reachable support after finitely many iterates. On that support, positive cost potentials and terminal-score equations are necessary and sufficient, with exact resolvent values. Unreachable recurrent blocks must be excluded rather than treated as physical nontermination. With finite algebraic primitive data and a compact semialgebraic decoder, the same conditions give a finite algebraic synthesis problem and an attained memory-constrained optimum. Sequence form, compact policy arguments, quantum decision processes, potential equations and real-algebraic elimination have established antecedents \cite{Stengel96,Balder89,Barry14,Ying18,Saldi25,Seidenberg54,Chatterjee15,Andriushchenko23}. Our statements preserve their scope and identify the particular common-model, all-initial-state, paid and finite-label certificates being combined; they do not claim those general methods as new.

Third, a natural controlled-measurement application yields an exact spectral resource law rather than another engineered symmetrization. Consider two projective detectors with failure projections $P,Q$, no additional control operation, unit call cost, and an arbitrary unknown initial state. When both projections are nonzero and $c=\|PQ\|<1$, the unrestricted-history optimum is $1+[2(1-c)]^{-1}$ and the one-live-label optimum is $1+(1-c)^{-1}$. Theorem~\ref{thm:two-label-exact} gives the exact two-label optimum, with a unique controller parameter and a transition at $c^2=(\sqrt5-1)/2$. Three fixed-start live labels attain the unrestricted value. The finite-dimensional two-subspace geometry is classical \cite{Halmos69}; the resource theorem proves the complete integer label frontier and its optimal strategies for the stated detector interface. It is distinct from first-detection models with one target subspace and intervening unitary evolution or optimized probing time \cite{Kessler21,DelVecchio26}.

The detector model also exposes a resource-sensitive minimax issue. A single explicitly derived pure preparation is least favorable against every full-history policy. But at two fixed-start live labels the robust value is strictly larger than the value obtained by optimizing a policy separately for each preparation. Corollary~\ref{cor:memory-minimax-gap} computes this gap exactly through the two-label frontier. Randomizing the initial counted label changes the interface and removes it. Thus one may neither insert a minimax exchange nor declare a persistent initial random seed free. This example does not invalidate a reference-based lower theorem; it explains why a sharp reference witness is an additional structural requirement in a nonconvex resource class.

The inherited acquired geometry and realizations remain part of the article. They include non-power sublevel comparisons, regular and mixed-dimensional mass laws, zero-loss atoms, singular task matrices, noncommuting quantum instruments, conditional-hazard implementation, calibrated causal composition and a Borel predictive quotient constructed on compact positive-evidence history patches. The complete preceding technical companion retains nonintervening directional Gaussian sensing, exact stopping resources, absolute-calibration occupation bounds and finite-description examples \cite{Companion}. The original ordered-measurement, two-stage reset, rank-constrained discrimination, variational-dual and spectral developments remain in their pinned independent realization sources, recorded in the preservation map. They are not silently deleted or relabelled as consequences of the new detector theorem.

These layers have different hypotheses. The transfer is valid for the declared broad prepared-kernel class with its reference and construction certificates. The finite-obstruction classification uses finite branching and a positive cost floor. Fixed-memory algebraic synthesis uses a finite time-homogeneous model family and explicit finite primitive data. The closed detector formulas concern two specified projectors and their physical call convention. Local acquired geometry does not by itself classify every experiment, and none of these theorems establishes the independent A1, A2 or unrestricted Theta-pipeline targets. Appendix~\ref{sec:current-scope} and the current technical supplement make the distinctions explicit. Throughout, program description, scratch, precision, calibration and execution time remain separately charged rather than inferred from the number of retained labels.
